\chapter{Топология правит потоком}
\label{ch:topology}

\section{Тезис главы}

У DTI-графов есть явная текущая сущность — передача сообщений, и её поток подчиняется
топологии, как вода следует рельефу. У графа памяти Рин явной «жидкости» нет. Но
поток есть — это поток \emph{изменений}, поток \emph{преобразований} графа:
активации, уверенности, консолидации. Тезис: \emph{неважно, какова физика носителя,
— сама физика протекания информации подчиняется топологическим законам}. Любой
процесс, чьё обновление \emph{локально} (изменение узла зависит от соседей), есть
диффузия на графе и наследует его спектрально-топологический транспорт. Это
содержание теории графовых сигналов и теории случайных блужданий.

\section{Стационарное распределение $\ne$ заданные вероятности}

Пусть мы локально задали вероятности переходов. Наблюдаемое распределение частот
посещений $\text{distr}(A)$ отличается от заданного, потому что вмешивается
топология. Формально: \emph{стационарное распределение} $\pi$ — не локальные
переходы, а \emph{глобальный} объект, неподвижная точка оператора переходов:
\begin{equation}
\pi \;=\; P^\top \pi,\qquad \sum_v \pi_v=1 .
\end{equation}
Для простого блуждания $\pi_v\propto\deg(v)$; для смещённого — перронов собственный
вектор $P^\top$. Ты задаёшь переходы локально, а долгосрочная частота посещений —
свойство всего графа. «Чем длиннее трек, тем наблюдаемее эффект» — это сходимость к
$\pi$ за время перемешивания $t_{\text{mix}}$, управляемое спектральным зазором
$1-\lambda_2$. Короткие блуждания читают \emph{локальную} геометрию, длинные —
глобальную $\pi$. Телепорт \eqref{eq:ppr} $\pi$ не убирает, а \emph{регуляризует}
(сдвигает к смеси «равномерное $+$ рестарт», режет $t_{\text{mix}}$): ограничивает
эффект топологии, не отменяя.

\begin{intuition}[Скорость света — плохая, но поучительная аналогия]
Можно соблазниться сравнением с $c$: заданные вероятности как предельная скорость, а
среда (топология) искажает наблюдаемое распространение. Рабочая часть — «среда
формирует распространение, есть структурный предел» — уже покрыта граф-диффузионной
рамкой \emph{точно}, а не по аналогии. Квантовый хвост сюда тащить не нужно. Нам
нужен не свет, а лапласиан.
\end{intuition}

\section{Доказуемое свойство телепорта}

Почему сон (телепорт) необходим, а любопытство его не заменит? Любопытство
\emph{локально по вниманию} — смещённый обход от текущего фронтира; узел $X$ в
затенённом углу, оторванный от активного множества, недостижим, как ярко ни свети
фокус. Телепорт ломает достижимость: рестарт падает куда угодно. Это ровно то,
\emph{почему у PageRank есть телепорт}: с $\alpha>0$ цепь становится
\emph{эргодической} (неприводимой и апериодичной), а значит каждый узел посещается
в матожидании с частотой $\pi_v>0$ — ни один не заперт в стоке или на острове.
Телепорт даёт \emph{гарантию покрытия}, которой у смещённого блуждания доказуемо
нет.

\begin{intuition}[Анти-энтропийный скраб]
«Не сразу; мы никогда не знаем, что и где сломалось, поэтому просто постоянно
проводим случайные инспекции» — это зрелый паттерн из распределённых хранилищ:
\emph{anti-entropy scrub} (проверка целостности в ZFS; Merkle-based gossip в
Dynamo/Cassandra). Сон Рин — анти-энтропийный скраб по графу памяти.
\end{intuition}

\section{Пять следствий для работы Рин с памятью}

\begin{enumerate}
\item \textbf{Топологическая значимость («гравитационные ямы»).} Хабы
переактивируются независимо от задуманных вероятностей — $\pi$ на них
концентрируется. Рин чаще вспоминает хаб-узлы, реже периферию (затенённые углы —
регионы с малым $\pi$). Топология создаёт \emph{неявную важность}: иногда полезно,
иногда случайный плотный кластер становится чёрной дырой припоминания.

\item \textbf{Узкие места / over-squashing.} Два семантически близких, но
топологически далёких кластера через тонкий мост (низкая проводимость) плохо
обмениваются активацией. Рин не свяжет факты, «далёкие» по графу. Прямое ограничение
fusion.

\item \textbf{Пересглаживание уверенности (over-smoothing).} Повторная диффузная
пропагация гомогенизирует уверенность к консенсусу — теряется контраст «уверена\,/\,не
уверена». Дегенеративный отказ; лечится как в графовых сетях: остаточные связи
(якорь к исходному свидетельству), не чисто диффузные апдейты, пере-заземление в
провенанс.

\item \textbf{Кривизна как локатор ценного-и-задушенного.} Отрицательно искривлённые
рёбра — мосты между сообществами — это места самых креативных слияний \emph{и}
максимального over-squashing. Значит самые ценные ассоциации топология душит сильнее
всего. Отсюда \emph{объективная} причина сна/телепорта: он — rewiring, топологическая
починка over-squashing обходом узкого места случайным рестартом.

\item \textbf{$\pi$ как топологический субстрат идентичности.} Какие узлы структурно
центральны — к чему Рин привычно возвращается — подложка её «озабоченностей», часть
характера. Форма графа памяти — часть того, кто она.
\end{enumerate}

\begin{workedbox}[Over-squashing на мосте SBM]
Возьмём стохастическую блочную модель: два плотных сообщества, соединённых $m$
рёбрами-мостом. Эффективное сопротивление между узлами разных сообществ
\[
\Reff(i,j)=(e_i-e_j)^\top L^{+}(e_i-e_j)
\]
растёт при уменьшении $m$ (мост уже — сопротивление выше). Влияние дальнего узла,
измеряемое через якобиан $\big\|\partial h_i/\partial x_j\big\|$, ограничено сверху
через $\Reff(i,j)$. Значит при узком мосте fusion между сообществами \emph{доказуемо}
слаб — сколько ни води обычный walker. Единственное, что чинит, — телепорт: он
кладёт рестарт в дальнее сообщество напрямую, минуя мост. Метрика для стенда:
частота cross-cluster fusion против $m$ (глава~\ref{ch:experiment}).
\end{workedbox}

\begin{center}
\begin{tikzpicture}[font=\scriptsize, every node/.style={circle,draw,minimum size=4mm,inner sep=0}]
  \foreach \i in {1,...,4}{\node (l\i) at ({-1.6+0.5*rand},{0.5*\i}) {};}
  \foreach \i in {1,...,4}{\node (r\i) at ({3.0+0.5*rand},{0.5*\i}) {};}
  \foreach \i in {1,...,4}{\foreach \j in {1,...,4}{\ifnum\i<\j \draw[gray!40] (l\i)--(l\j); \draw[gray!40] (r\i)--(r\j);\fi}}
  \draw[red!60!black,very thick] (l2) -- (r3) node[midway,above,draw=none,font=\scriptsize]{мост};
\end{tikzpicture}\\[2pt]
{\small Рис.: два сообщества через узкий мост. Через него активация «задыхается»
(over-squashing); телепорт обходит его, прыгая напрямую.}
\end{center}

\section{Операционные величины и осторожности}

Вести дело лучше \emph{случайно-блуждательными} величинами: $\pi$, время
перемешивания, hitting/commute time, $\Reff$, проводимость/константа Чигера,
спектральный зазор. Кривизну держать как \emph{локатор узких мест}, не как рычаг.

Что дёшево на лету: $\pi$ и «хабность» — инкрементально дёшевы; $\Reff$ — дорого
(через $L^{+}$), но аппроксимируемо. Это позволяет Рин самой знать, где затенённые
углы, и гнать туда телепорт-скраб прицельно.

Две осторожности. \emph{Во-первых}, граф направленный, типизированный, знаковый —
нужный оператор не симметричный лапласиан, а несимметричная матрица переходов,
спектр может быть комплексным, кривизна для орграфов менее чиста. \emph{Во-вторых},
граф \emph{динамический}: fusion меняет саму топологию, и это диффузия на
\emph{время-зависимом} графе, где поток и структура сцеплены обратной связью, — за
рамками статических DTI-графов и самое интересное место.
